\ifdefined\XeTeXversion\else\pdfoutput=1\fi
\documentclass[10pt]{nvidiaextabstract}
\usepackage{xurl}
\ifdefined\XeTeXversion\microtypesetup{tracking=false}\fi
\venue{From Human Videos to Dexterous Robot Capability (CoRL 2026)}
\submissionID{longbot -- Track 1}
\title{Auditable Monocular Human--Object Reconstruction\\with Evidence-Constrained Repairs}
\shorttitle{longbot: V2D Track 1}
\author[1,2]{Ismail Syed}
\affil[1]{University of Texas at Austin}
\affil[2]{\href{https://www.longhorizon.bot}{www.longhorizon.bot}}
\begin{abstract}
We reconstruct human geometry, object meshes and continuous object motion from
30 monocular V2D demonstrations using pretrained visual models, optimization
and targeted repairs. Across thirteen submitted variants, the retained T1-11
candidate reduces object Chamfer distance from 25.60394 to 24.21106 cm and object
acceleration error from 1.74271 to 0.49967 cm/frame$^2$ relative to our initial
submission. Human geometry improves slightly, while penetration worsens.
We contribute an auditable sequence of positive and negative experiments,
source-constrained repair mechanisms and reusable interfaces. The results
support cautious geometric/temporal correction, not universal improvement
from additional optimization or a different pretrained body estimator.
\end{abstract}
\begin{document}
\maketitle
\section{Introduction}
Monocular human--object reconstruction must resolve identity, geometry, scale,
occlusion and motion from ambiguous observations. V2D Track 1~\cite{v2d}
evaluates both reconstruction accuracy and physical plausibility against
withheld multi-view estimates. The output comprises MHR human parameters,
object geometry and per-frame object poses, exported on the required scored
frame grid. Track 1 and Track 2 episode indices describe different sequences;
our Track 1 reconstruction does not use Track 2 assets or trajectories.

We build on a CARI4D-style workflow~\cite{cari4d} and evaluate targeted
changes against an immutable retained baseline. Our contribution is the
repair, validation and provenance workflow, rather than a new pretrained
vision model. This report designates the complete T1-11 file; later T1-12 and
T1-13 trials did not justify replacing it.

\section{Approach}
\textbf{Initialization and refinement.}
The reconstruction stack combines SAM2 segmentation~\cite{sam2}, monocular
depth, SAM 3D body/object estimates~\cite{sam3d}, object registration/tracking
and joint human--object refinement. We investigated CoCoNet and the released
iterative refinement path; later diagnostic pipelines also used MoGe-2 and
FoundationPose. Their experiments are distinguished from components that
actually produced the retained submission. Human parameters, triangle meshes,
object rotations/translations and positive object scale are exported with
explicit coordinate conventions and timestamps.

\textbf{Targeted repairs.}
We investigated geometry scale/proportion, orientation ambiguity, temporal
noise, support consistency and person/object selection. Candidate transforms
are bounded by source-image evidence and preserve unmodified intervals.
Object-family experiments test shape-appropriate hypotheses instead of
applying the same correction to all objects. Accepted changes are layered on
the retained baseline with per-episode identity and provenance checks.
Occluded frames retain a continuous trajectory as required by evaluation,
but inferred motion is not treated as an independently observed pose.

\textbf{Evaluation discipline.}
The challenge's first-frame human Sim(3) alignment affects both human and
object evaluation. Changing human scale or placement can therefore alter
object scores even when object coordinates are unchanged. We compare human,
object and physical-plausibility metrics jointly. Visual silhouette agreement,
source depth diagnostics and official scores are separate evidence streams.
Image-gated repair hypotheses are generated without hidden reference inputs;
submission selection uses observed benchmark scores and may overfit this small
set. Rejected hypotheses and producing source snapshots remain available.

\section{Experiments \& Results}
We submitted thirteen complete variants spanning initialization/refinement,
geometry and temporal corrections, object-family repairs and alternative
human estimates. Table~\ref{tab:results} gives exact scores of the initial and
retained 30-episode files. All metrics are lower-is-better. Chamfer distances
and penetration are in centimetres; acceleration discrepancies are in
cm/frame$^2$ under the metric's discrete-frame convention.

\begin{table}[t]
\centering\small\setlength{\tabcolsep}{4pt}
\begin{tabular}{lrr}
\toprule
Metric & Initial T1-01 & Retained T1-11\\
\midrule
CD-H & 15.53707 & 15.47123\\
CD-O & 25.60394 & 24.21106\\
ACC-H & 1.38088 & .27484\\
ACC-O & 1.74271 & .49967\\
PEN & .00430 & .00485\\
\bottomrule
\end{tabular}
\caption{Kaggle scores of two exact files. The retained file covers all 30
episodes and 9,877 scored frames. These are not independent-seed averages.}
\label{tab:results}
\end{table}

The retained result reduces CD-O by 5.44\% and ACC-O by 71.33\%.
CD-H changes little and PEN becomes worse: smoothing and improved pose
agreement do not ensure better contacts. An early geometry/temporal variant
reduced acceleration error but worsened CD-O to 32.65560 cm; selective
rollback recovered object accuracy. Subsequent object-family work decreased
CD-O from 24.82940 through 24.31892 to 24.21106 cm. The final retained
iron-family change affects Episode 8 while preserving all human values and
unselected object values. Later desk/proportion submissions yield CD-O
24.45089 and 24.21206 cm and remain diagnostic alternatives.

\textbf{Negative studies.}
More refinement was not uniformly beneficial, motivating bounded comparisons
and checkpoint reuse. Alternative human pipelines were tested, including
FlowHMR~\cite{flowhmr} and GVHMR diagnostics. A wrong-person FlowHMR case
improved held-out silhouette IoU from 0.3303 to 0.7222 after selection repair,
but remained below the retained baseline's 0.8029. Across the all-30 image
screen, mean baseline IoU was 0.7855 versus 0.7085 for the corrected official
FlowHMR branch. Representation-conversion error, camera-depth ambiguities
and body identity must be evaluated separately; a successful conversion
check is not ground-truth reconstruction accuracy. These unselected studies
are documented but are not claimed as contributors to the retained score.

\section{Conclusion}
Evidence-constrained, object-specific corrections improve the aggregate
reconstruction, but accuracy, temporal fidelity and contact plausibility can
conflict. \textbf{Economic viability} is addressed through pretrained components,
selective recomputation and checkpointed experiments on laptop/L4 hardware;
total production cost savings have not been measured.
The \textbf{data flywheel} stores diagnosed failures, verified repairs and
regression cases for reuse, while acknowledging continued manual review.
\textbf{Complexity reduction} comes from explicit geometry/time/coordinate
contracts and modular reconstruction/repair interfaces rather than new capture
hardware. \textbf{Innovation} is treating repairs as testable hypotheses with
rollback and preserving negative results instead of assuming every refinement
or model replacement improves reconstruction.

Future work should prioritize metric-depth observability, contact and
cross-sequence validation. This is a challenge-scale study, not a demonstrated
autonomous production system. Code and provenance:
\url{https://github.com/isyed/v2d-reconstruction}; reusable modules:
\url{https://github.com/isyed/real2sim2real}. Exact submission IDs and the
producing commit accompany the report; packaged-output verification and
raw-video reruns are reported separately.
\bibliographystyle{plainnat}
\bibliography{references}
\end{document}
